\section{Conclusion}
We study on-policy distillation for language models and identify a key limitation of the reverse-KL objective:
its mode-seeking nature can collapse generation diversity and destabilize learning at token positions where the teacher distribution has high entropy.
% its mode-seeking behavior can collapse generation diversity and produce unstable learning signals at token positions where the teacher distribution has high entropy. 
To address this, we introduce \method, which adapts the on-policy distillation objective and selectively applies forward-KL for high-entropy teacher tokens. Empirically, \method~better preserves the teacher-like distribution while retaining on-policy training efficiency, yielding consistent gains over standard on-policy distillation on six mathematical reasoning benchmarks.
% . Across six mathematical reasoning benchmarks, this leads to consistent improvements over standard on-policy distillation. 
Overall, our results demonstrate that explicitly modeling teacher uncertainty is essential for stable, diverse, and effective knowledge transfer.

\section*{Acknowledgements}
This work was supported by Institute for Information \& communications Technology Planning \& Evaluation(IITP) grant funded by the Korea government(MSIT) (RS-2019-II190075, Artificial Intelligence Graduate School Program(KAIST)), the MSIT(Ministry of Science, ICT), Korea, under the Global Research Support Program in the Digital Field program(RS-2024-00436680) supervised by the IITP(Institute for Information \& Communications Technology Planning \& Evaluation), and the Institute of Information \& Communications Technology Planning \& Evaluation(IITP) grant funded by the Korea government(MSIT) (RS-2025-02304967, AI Star Fellowship(KAIST))